% !TEX program = xelatex
% 暨南大学本科生课程论文 — SIS∞ 求解框架
% 编译：xelatex 课程论文_SIS无穷范数求解.tex（需 ctex 宏包）
\documentclass[12pt,a4paper]{ctexart}
\usepackage{amsmath,amssymb,amsfonts}
\usepackage{graphicx}
\usepackage{booktabs}
\usepackage{hyperref}
\usepackage{geometry}
\usepackage{indentfirst}
\usepackage{caption}
\usepackage{subcaption}
\usepackage{listings}
\usepackage{xcolor}
\geometry{left=2.5cm,right=2.5cm,top=2.5cm,bottom=2.5cm}

\lstset{
  basicstyle=\small\ttfamily,
  breaklines=true,
  frame=single,
  numbers=left,
  numberstyle=\tiny,
  keywordstyle=\color{blue},
  commentstyle=\color{gray},
}

% 封面信息（请自行填写）
\newcommand{\papertitle}{面向 SIS$_{\infty}$ 的双空间启发式搜索与 CP-SAT 收尾：\\2026 密码数学挑战赛赛题一求解实践}
\newcommand{\college}{【请填写学院】}
\newcommand{\major}{【请填写专业】}
\newcommand{\studentname}{【请填写姓名】}
\newcommand{\studentid}{【请填写学号】}
\newcommand{\coursename}{【请填写课程名称】}
\newcommand{\advisor}{【请填写指导教师】}
\newcommand{\submitdate}{2026年6月}

\begin{document}

% ========== 封面 ==========
\begin{titlepage}
  \centering
  \vspace*{1.5cm}
  {\LARGE\bfseries 暨\hspace{0.5em}南\hspace{0.5em}大\hspace{0.5em}学\par}
  \vspace{0.8cm}
  {\Large 本科生课程论文\par}
  \vspace{2.5cm}
  {\large\bfseries 论文题目：\par}
  \vspace{0.5cm}
  {\Large\bfseries \papertitle\par}
  \vfill
  \begin{tabular}{rl}
    学\hspace{2em}院： & \college \\[0.6em]
    专\hspace{2em}业： & \major \\[0.6em]
    学生姓名： & \studentname \\[0.6em]
    学\hspace{2em}号： & \studentid \\[0.6em]
    课程名称： & \coursename \\[0.6em]
    指导教师： & \advisor \\
  \end{tabular}
  \vfill
  {\large \submitdate\par}
\end{titlepage}

% ========== 摘要页 ==========
\begin{center}
  {\Large\bfseries \papertitle}\\[1em]
\end{center}

\noindent\textbf{【摘要】}
2026 年密码数学挑战赛赛题一要求在模 $q$ 整数矩阵 $A$ 上求解短整数向量 $(u,v)$，满足同余约束 $Av+u\equiv t\pmod q$ 且 $\|u\|_{\infty},\|v\|_{\infty}\le\gamma$。该问题在欧氏范数格约化（BKZ/LLL）与无穷范数目标之间存在本质失配。本文提出并实现了一套\textbf{双空间残差优先（Dual-Space Residual-First）}求解框架：在 $v\in[-\gamma,\gamma]^m$ 上最小化中心化残差 $r(v)=\mathrm{Center}(t-Av)$ 的 $L_{\infty}$ 违规度，结合 fpylll 格种子、模 $q$ 核游走、Wagner 子系统碰撞、行级 CP-SAT 修补与全维/分块/分层整数规划收尾。针对赛题三类（齐次 SVP、非齐次 CVP、欧氏下界特殊题）设计差异化参数调度。在官方实例上的系统实验表明：启发式搜索可将第一类四题的最坏坐标违规压至 $44\sim46$；引入 ortools 全维 $v$ 的 CP-SAT 收尾后，题 1/3/6/9 分别可达 $\|u\|_{\infty}=42,43,44,45$（$\gamma=15$），相对 batch 基线平均降低约 2 个 $L_{\infty}$ 量级，但仍距可行域差约 27--30。分块 ILP 与小子格 BKZ+局部搜索在已平台实例上未进一步突破。本文总结了算法组件贡献、瓶颈诊断与后续研究方向。

\vspace{0.8em}
\noindent\textbf{【关键词】}短整数解；无穷范数；格约化；双空间搜索；CP-SAT；密码数学挑战赛

\newpage
\tableofcontents
\newpage

% ========== 正文 ==========
\section{引言}

\subsection{研究背景}

短整数解问题（Short Integer Solution, SIS）是现代格密码学的核心困难假设之一。给定模 $q$ 矩阵 $A\in\mathbb{Z}_q^{n\times m}$ 与目标向量 $t\in\mathbb{Z}_q^n$，经典 SIS 要求寻找非零短向量 $x$ 使得 $Ax\equiv 0\pmod q$。2026 年密码数学挑战赛赛题一将其扩展为\textbf{无穷范数}版本（记为 SIS$_{\infty}$）：求 $u\in\mathbb{Z}^n,v\in\mathbb{Z}^m$ 使得
\begin{equation}
  Av + u \equiv t \pmod q,\qquad
  \|u\|_{\infty}\le\gamma,\quad \|v\|_{\infty}\le\gamma.
  \label{eq:sisinf}
\end{equation}
官方参数规模约为 $n=m=q\in\{100,120,140,160\}$，$\gamma\in\{15,16,17,18\}$。与以欧氏范数为目标的 BKZ 算法不同，本题直接优化 Chebyshev（$L_{\infty}$）范数，高维下二者可相差 $\Omega(\sqrt{n})$ 量级，构成主要算法挑战。

\subsection{赛题结构与本文工作}

赛题按数学结构分为三类（见表~\ref{tab:taxonomy}）：第一类（题 1/3/6/9）为齐次情形 $t\equiv 0$，侧重 SVP 与核空间游走；第二类（题 2/4/7/10）为非齐次 $t\not\equiv 0$，侧重 CVP 嵌入与模拉回种子；第三类（题 5/8）在式\eqref{eq:sisinf} 之外另需 $\|u\|_2^2+\|v\|_2^2\ge q^2$。

\begin{table}[htbp]
  \centering
  \caption{赛题三类划分与求解侧重}
  \label{tab:taxonomy}
  \begin{tabular}{clll}
    \toprule
    类别 & 题号 & 特征 & 求解侧重 \\
    \midrule
    一 & 1,3,6,9 & 齐次，禁平凡零解 & BKZ/LLL 种子、核游走、$u$ 优先算子 \\
    二 & 2,4,7,10 & 非齐次 & CVP 提升、模拉回、弱化 BKZ \\
    三 & 5,8 & $L_{\infty}$+欧氏下界 & 熵/欧氏权重、分层 ILP、欧氏抛光 \\
    \bottomrule
  \end{tabular}
\end{table}

本文在开源实现 \texttt{sisinf\_challenge2026} 中完成了：（1）统一数学建模与三类题参数调度；（2）双空间启发式主搜索器；（3）可选 fpylll 真格约化种子管线；（4）ortools CP-SAT 多模式收尾（全维/分块/分层）；（5）十题批量流水线 \texttt{run\_all\_pipeline.py}；（6）系统消融与长跑实验。下文第 2 节阐述原理与方案，第 3 节分析实现与实验，第 4 节总结结论与展望。

\section{原理或方案设计}

\subsection{残差优先等价变换}

定义对称中心化映射 $\mathrm{Center}:\mathbb{Z}_q\to\mathbb{Z}\cap[-q/2,q/2]$。对固定 $v$，令
\begin{equation}
  r(v)=\mathrm{Center}(t-Av\bmod q),\qquad u=r(v).
\end{equation}
则式\eqref{eq:sisinf} 的可行性等价于：存在 $v\in[-\gamma,\gamma]^m$ 使得 $\|r(v)\|_{\infty}\le\gamma$ 且 $\|v\|_{\infty}\le\gamma$。该变换将双变量 $(u,v)$ 搜索化为\textbf{单变量 $v$ 的盒约束优化}，且 $u$ 由残差闭式给出，避免同余与 $L_{\infty}$ 的耦合震荡。

\subsection{分层离散目标（Lexicographic Score）}

定义对 $(r,v)$ 的三级字典序目标（与程序 \texttt{objective\_uv} 一致）：
\begin{align*}
  \text{violations} &= \#\{i:|r_i|>\gamma\}+\#\{j:|v_j|>\gamma\},\\
  \text{overflow\_sum} &= \sum_i\max(0,|r_i|-\gamma)+\sum_j\max(0,|v_j|-\gamma),\\
  \text{max\_overflow} &= \max\bigl\{\max_i(|r_i|-\gamma),\max_j(|v_j|-\gamma),0\bigr\}.
\end{align*}
搜索接受规则默认采用单调改进（禁止上坡模拟退火），以在有限离散状态空间上稳定收敛至局部最优。

\subsection{Dual-Space 候选生成}

\textbf{双空间}指：在 $v$ 空间生成多样初值，在残差空间评估与修复。

\paragraph{格种子（Lattice Seeds）}
构造 Ajtai 型格基：对 $B\in\mathbb{Z}^{(n+m)\times(n+m)}$ 前 $n$ 行嵌入 $qI_n$，后 $m$ 列对应 $-A$ 的列与单位块。使用 fpylll 做 LLL 预处理；当维数 $n+m\le 160$ 时进一步运行多 block\_size 的 BKZ tour（Li--Nguyen 式弱 oracle 多种子策略）。约化基向量后 $m$ 坐标裁剪为 $v$ 种子。当 $n+m>160$（如题 6/9）时仅 LLL，避免全维 BKZ 数小时级耗时。

\paragraph{Wagner 子系统种子}
对 worst-$u$ 行对应的列子块，采用 $4\times4$  meet-in-the-middle 碰撞与低维暴力后备，生成满足局部同余的小 $v$ 补丁（非完整 Wagner--Gauss 次指数算法，但工程上可快速扩充种子池）。

\paragraph{模拉回与非齐次 CVP 种子}
对 $\phi\in\{\mathrm{center}(t),\mathrm{sign}(t),\ldots\}$，构造 $v\propto -A^\top\phi$ 并裁剪至 $[-\gamma,\gamma]^m$；非齐次题另用最小二乘近似 $Av\approx t+qk$ 的 CVP lift 变体。

\paragraph{模 $q$ 核游走（Kernel Walk）}
预计算 $A v\equiv 0\pmod q$ 的核基 $K$，在 $v\leftarrow v+Kc$ 上随机游走（$c$ 小系数），在 $v$ 超界时协调降低 $L_{\infty}$ 违规。

\subsection{$u$ 优先局部算子（第一类强化）}

实验表明第一类题瓶颈在 $\|u\|_{\infty}$ 而非 $\|v\|_{\infty}$。针对性引入：
\begin{itemize}
  \item \textbf{ViolationLS / 分层行 LS}：对溢出最严重的坐标行做最小二乘投影；
  \item \textbf{u\_row\_snap}：在最坏 $u$ 行关联列上枚举 $\Delta v_j$；
  \item \textbf{CP-SAT 周期修补}：对 worst 行/强耦合列子集建立有界 $\Delta v_j$ 整数模型，最小化行溢出代理；
  \item \textbf{块 CP-SAT}：同时优化多行溢出上界。
\end{itemize}
消融实验（题 1，u 优先 B 档相对 A 档）可将 $\|u\|_{\infty}$ 从 46 降至 45 量级。

\subsection{CP-SAT 全维 $v$ 收尾（ILP Finish）}

当启发式平台化（如 $\|u\|_{\infty}\approx 44$）时，引入 ortools CP-SAT 对\textbf{全部} $v_j\in[-\gamma,\gamma]$ 建模。对每行 $i$ 引入整数 $k_i$ 实现中心化：
\[
  |u_i| = \bigl|\mathrm{center}(t_i-(Av)_i)\bigr| \le \gamma + \mathrm{over}_i,\quad \min \max_i \mathrm{over}_i.
\]
实现三种模式：
\begin{enumerate}
  \item \textbf{full}：一次性全维优化（默认用于类一、类二）；
  \item \textbf{chunk}：每轮仅释放与 worst-$u$ 强耦合的 40 列，12 轮滑动覆盖；
  \item \textbf{lex}：三阶段——(1) 最小化 $\max\mathrm{over}$；(2) 固定上界后最小化 $\sum\mathrm{over}$；(3) 减少顶在平台上的行数（默认用于类三）。
\end{enumerate}

\paragraph{小子格 BKZ 种子（Sub-lattice BKZ）}
在 worst-$u$ 对应的 $40\times40$ 子块上建 80 维 Ajtai 格并强制真 BKZ，将短向量嵌入全维 $v$ 后短程局部搜索。实验表明在题 1 上常使 $\|u\|_{\infty}$ 从 42 回退至 45，\textbf{默认关闭}。

\paragraph{第三类欧氏抛光}
在 $L_{\infty}$ ILP 后若 $\|u\|_2^2+\|v\|_2^2<q^2$，以高 \texttt{euclid\_weight} 与熵权重运行短程局部搜索，兼顾欧氏下界。

\subsection{三类题参数调度}

\texttt{apply\_sis\_class\_defaults} 按类别覆盖 \texttt{SearchConfig}：类一启用 BKZ($\beta\approx28$)、Wagner、$u$ 优先 CP；类二关闭 BKZ、增强 CVP/pull；类三提高欧氏与熵权重、缩短核游走启动进度。批量流水线 \texttt{run\_all\_pipeline.py} 对每题执行「batch $\times$6 轮 $\to$ ILP 3600s $\to$ 汇总报告」。

\section{程序实现或算法分析}

\subsection{系统架构}

核心模块如表~\ref{tab:modules}。

\begin{table}[htbp]
  \centering
  \caption{主要程序模块}
  \label{tab:modules}
  \small
  \begin{tabular}{ll}
    \toprule
    文件 & 功能 \\
    \midrule
    \texttt{solve\_sisinf.py} & 主求解器、\texttt{local\_search\_one}、CP-SAT 修补 \\
    \texttt{lattice\_bkz.py} & Ajtai 格、fpylll BKZ/LLL、启发式回退 \\
    \texttt{sis\_advanced\_u\_ops.py} & Wagner 种子、ViolationLS、离散高斯种子 \\
    \texttt{modq\_kernel.py} & 模 $q$ 右核基 \\
    \texttt{sis\_finish\_ops.py} & 全维/分块/分层 ILP \\
    \texttt{finish\_core.py} & 收尾编排、类三欧氏抛光 \\
    \texttt{run\_class\_batch.py} & 按类批量、保存 \texttt{problem*\_best.json} \\
    \texttt{run\_all\_pipeline.py} & 十题统一流水线 \\
    \bottomrule
  \end{tabular}
\end{table}

\subsection{复杂度与工程考量}

\begin{itemize}
  \item 单坐标移动：利用 $r\leftarrow \mathrm{Center}(r-\Delta v_j\cdot A_{:,j})$，摊还 $O(n)$；
  \item 全维 ILP：变量 $O(m)$，每行引入 $k_i$ 与 $|\cdot|$ 辅助变量，$m=100\sim160$ 时 CP-SAT 1 小时为可接受收尾预算；
  \item 维数 $n+m>160$：BKZ 退化为 LLL only；题 9（160 维）ILP 改善弱于题 1（100 维）。
\end{itemize}

\subsection{实验设置}

\begin{itemize}
  \item 环境：Linux，Python 3，numpy；fpylll（conda-forge）；ortools（独立 venv，避免 protobuf 冲突）；
  \item 批量：每题 6 轮渐进配置（\texttt{\_cfg\_for\_round}），round$\ge 6$ 启用 aggressive BKZ/CP 窗口；
  \item ILP：默认 3600s，4 workers；sub-BKZ 默认跳过。
\end{itemize}

\subsection{主要实验结果}

\subsubsection{第一类四题（$\gamma=15$，$t\equiv 0$）}

表~\ref{tab:class1} 汇总当前最优进展（\texttt{inf\_u} 即 $\|u\|_{\infty}$）。

\begin{table}[htbp]
  \centering
  \caption{第一类题启发式 batch 与 full ILP 收尾对比}
  \label{tab:class1}
  \begin{tabular}{ccccc}
    \toprule
    题号 & $n\times m$ & batch 最优 $\|u\|_{\infty}$ & ILP 后 $\|u\|_{\infty}$ & $\|v\|_{\infty}$ \\
    \midrule
    1 & $100\times100$ & 44 & \textbf{42} & 15 \\
    3 & $120\times120$ & 45 & \textbf{43} & 15 \\
    6 & $140\times140$ & 45 & \textbf{44} & 15 \\
    9 & $160\times160$ & 45 & 45 & 15 \\
    \bottomrule
  \end{tabular}
\end{table}

\paragraph{观察}
\begin{enumerate}
  \item 同余约束始终满足（\texttt{congruence\_ok=1}）；
  \item $\|v\|_{\infty}$ 在 batch 后多已贴界 $\gamma=15$，瓶颈集中在 $u=r(v)$；
  \item 全维 ILP 首轮通常降低约 2 个 $L_{\infty}$ 单位；题 1/3 在 7200s 二次 ILP 上\textbf{平台化}（42/43 不变）；
  \item 题 1 上 60 行 $u$ 违规、$v$ 96 维非零，说明非稀疏坏行，行级 snap 覆盖不足；
  \item \texttt{ilp\_max\_over=27} 与 $\|u\|_{\infty}=42=15+27$ 一致。
\end{enumerate}

\subsubsection{分块 ILP 与小子格 BKZ 消融}

\begin{table}[htbp]
  \centering
  \caption{题 1/3 上 chunk ILP（12 轮$\times$约 300s）结果}
  \label{tab:chunk}
  \begin{tabular}{cccc}
    \toprule
    题号 & incumbent $\|u\|_{\infty}$ & chunk 后 & 12 轮均 \texttt{improved}? \\
    \midrule
    1 & 42 & 42 & 否 \\
    3 & 43 & 43 & 否 \\
    \bottomrule
  \end{tabular}
\end{table}

小子格 BKZ+LS 在题 1 上将 $\|u\|_{\infty}$ 从 42 恶化至 45；heavy CP（\texttt{cp\_aggressive\_row\_k=60}）运行约 13h 仍停留在 42。

\subsubsection{u 优先参数消融（题 1，quick）}

B 档（默认类一配置）相对 A 档：$\|u\|_{\infty}$ 45 vs 46，\texttt{max\_overflow} 30 vs 31，验证 $u$ 优先算子有效但未达可行。

\subsubsection{格后端}

\texttt{lattice\_backend=fpylll}，\texttt{lattice\_seed\_count=40}（含 Wagner+BKZ 合并）。维数 200 时 fpylll 路径为 LLL only，仍优于纯启发式单位列种子。

\subsection{瓶颈分析}

当前框架在官方第一类实例上形成稳定\textbf{$L_{\infty}$ 平台}（42--45），距可行 $\gamma=15$ 差约 27--30。原因包括：
\begin{itemize}
  \item 欧氏 BKZ 短向量与 $L_{\infty}$ 目标失配；
  \item 全维 CP-SAT 在 $m\approx100$ 时 1h 内仅能保证局部改进；
  \item $v$ 近满支撑时，分块 ILP 无法协调全局 $u$ 下降；
  \item 赛题实例可能位于当前启发式+CP 能力边界之外，需更强非启发式攻击（高 $\beta$ 子格、Kannan 嵌入、完整 Wagner 等）。
\end{itemize}

\section{总结}

本文针对 2026 密码数学挑战赛 SIS$_{\infty}$ 赛题，提出了\textbf{双空间残差优先}求解框架，将同余短向量问题转化为盒约束下 $L_{\infty}$ 残差最小化；结合 fpylll 格种子、模 $q$ 核游走、Wagner 子系统碰撞、$u$ 优先算子与 ortools CP-SAT 多模式收尾，并针对三类题设计差异化调度与十题批量流水线。

\textbf{主要创新点}包括：（1）残差闭式消元与 lex 离散目标的一致性实现；（2）真/假 BKZ 分层策略（全维 LLL vs 小子格强制 BKZ）；（3）全维/分块/分层 $v$ 的 CP-SAT $L_{\infty}$ 收尾与第三类欧氏抛光；（4）可复现的消融与 pipeline 工程。

\textbf{实验结论}：第一类四题在 ILP 后最优 $\|u\|_{\infty}$ 为 42--45，均未达到 $\gamma=15$；ILP 相对启发式平均改善约 2；chunk 与 sub-BKZ 在已平台实例上无效或有害。第二类、第三类十题流水线已就绪，完整长跑结果待补充。

\textbf{展望}：Kannan CVP 嵌入种子、更高 $\beta$ 受限子格 BKZ、lexicographic 多轮 ILP 热启动、SAT/ILP 子问题分解，以及结合题面特殊结构的理论剪枝。

\section*{致谢}
感谢课程指导教师在格密码与算法设计方面的指导；感谢开源社区 fpylll、ortools 与 numpy 生态。

\begin{thebibliography}{99}
\bibitem{ajtai1996}
M. Ajtai. Generating hard instances of lattice problems. \textit{STOC}, 1996.

\bibitem{micciancio2002}
D. Micciancio, O. Regev. Worst-case to average-case reductions based on Gaussian measures. \textit{FOCS}, 2004.

\bibitem{chen2011}
Y. Chen, P. Nguyen. BKZ 2.0: Better lattice security estimates. \textit{ASIACRYPT}, 2011.

\bibitem{wagner2002}
D. Wagner. A generalized birthday problem. \textit{CRYPTO}, 2002.

\bibitem{linfty}
本文赛题组. 2026 密码数学挑战赛赛题一：无穷范数下的短整数解问题（SIS$_{\infty}$）.

\bibitem{ortools}
Google OR-Tools. CP-SAT Solver Documentation. \url{https://developers.google.com/optimization}

\bibitem{fpylll}
fpylll: Python interface for LLL and BKZ. \url{https://github.com/fplll/fpylll}
\end{thebibliography}

\appendix
\section{附录：主要命令与目录结构}

\subsection{环境}
\begin{lstlisting}[language=bash]
conda install -c conda-forge fpylll
pip install -r requirements.txt
pip install -r requirements-ortools.txt  # 建议独立 venv
python scripts/check_fpylll.py
python scripts/check_ilp_finish.py
\end{lstlisting}

\subsection{单题求解}
\begin{lstlisting}[language=bash]
python scripts/solve_sisinf.py \
  --input saiti1/sis_inf_problems_json/problem1.json \
  --output results/solutions.json
\end{lstlisting}

\subsection{批量与收尾}
\begin{lstlisting}[language=bash]
# 第一类批量
python scripts/run_class_batch.py --class 1 --max-rounds 6

# ILP 收尾（auto：类一/二 full，类三 lex）
python scripts/finish_from_best.py \
  --instance saiti1/sis_inf_problems_json/problem1.json \
  --incumbent results/class1/problem1_best.json \
  --output results/p1_finish.json \
  --ilp-mode auto --ilp-time-limit 3600

# 十题流水线
python scripts/run_all_pipeline.py --batch-rounds 6 --ilp-time-limit 3600
\end{lstlisting}

\subsection{结果文件}
\begin{itemize}
  \item \texttt{results/classN/problem*\_best.json}：batch 最优 incumbent；
  \item \texttt{results/pipeline/all\_pipeline\_report.json}：十题汇总；
  \item \texttt{meta.lattice\_backend}：\texttt{fpylll} / \texttt{heuristic}。
\end{itemize}

\end{document}
